\begin{textsample}{POS Dim 1 – vem\_ed\_05 – Score 21.00 – vem\_ed\_05/t018.txt}  \label{ex:f1_pos_102}
QUEM \textbf{É} QUEM Brenda Ivonne Garcia Portillo Brenda Ivonne Garcia Portillo \textbf{é} diretora de projetos de internacionalização da Universidade de Monterrey (UDEM), uma das principais instituições no mundo em Intercâmbios Virtuais.
Ela relata os \textbf{pontos} altos dessa história de sucesso.
A UDEM \textbf{começou} a fazer COIL (Collaborative Online International Learning, uma \textbf{forma} de \textbf{intercâmbio} virtual) há cerca de 10 anos.
Fomos os primeiros no México a colaborar com o SUNY COIL Center [ EUA ].
Posteriormente nos afiliamos ao consórcio Global Partners for Education, liderado pela East Carolina University [ EUA ].
Decidimos replicar a metodologia COIL com os parceiros bilaterais.
Isso tomou tempo: divulgar o método, \textbf{estabelecer} \textbf{processos}, participar em iniciativas como EVALUATE da Unicollaboration.
Criamos oficinas de capacitação para docentes, além de publicações.
Assim, a UDEM se \textbf{tornou} \textbf{referência} internacional, presente no comitê organizador da conferência IVEC, membro fundador da LATAM COIL, entre outras \textbf{ações}.
As atividades da abordagem COIL devem se ligar aos objetivos de aprendizagem e incluir o \textbf{elemento} internacional.
As atividades \textbf{são} o meio para cumprir tais objetivos e, para enriquecê-las, o design de projeto deve selecionar as melhores ferramentas e plataformas tecnológicas.
Os projetos COIL se relacionam com a teoria de aprendizagem experiencial de Kolb, pois, ao participar em projetos \textbf{colaborativos} internacionais, os estudantes aprendem por meio da experiência.
Desenvolvem \textbf{empatia}, tolerância, respeito.
Trabalham com equipes internacionais, relacionam-se com indivíduos de outras culturas, adaptam-se à \textbf{mudança} e tornam-se mais flexíveis. \textbf{São} as \textbf{competências} interculturais, imprescindíveis no mundo globalizado.
Em um projeto \textbf{bem-sucedido}, \textbf{todos} contribuem para atingir as metas \textbf{compartilhadas}, encontrando o equilíbrio e a sinergia entre o acadêmico, o formativo e o administrativo.
A chave para o êxito da UDEM \textbf{é} o compromisso com as \textbf{estratégias} de internacionalização integradas à prática \textbf{institucional}.
Ter um plano estratégico definido \textbf{permite} \textbf{saber} aonde vamos e como atingir os objetivos.
Em 2020 praticamente duplicamos a participação em projetos COIL, com 1473 estudantes envolvidos em 65 grupos, sob orientação de 45 professores participantes e 53 universidades parceiras em 21 países.
Na Direção de Programas Internacionais da UDEM, existe uma Gerência de Projetos de Internacionalização em Casa, \textbf{que}, entre outras atividades, tais como Programa de Formação em Competências Interculturais e Feira Internacional, coordena os projetos COIL.
Somos apenas três na equipe e o trabalho tem \textbf{sido} intenso.
Esperamos crescer para dar um seguimento mais adequado às \textbf{estratégias} de internacionalização.

% matched lemmas: ação, bem-sucedido, colaborativo, começar, compartilhar, competência, elemento, empatia, estabelecer, estratégia, forma, institucional, intercâmbio, mudança, permitir, ponto, processo, que, referência, saber, ser, todo, tornar
\end{textsample}
